\section{结论与展望}
\label{sec:conclusion}

本文针对大模型介入摄像头视频异常检测后暴露出的三个真实问题，提出了一个统一框架 \evivad{}。我们的出发点是一个判断：在免训练范式已经证明「大模型可用于 VAD」之后，继续在检测分数上刷分的边际收益正在迅速下降，而\emph{解释是否可被核验}与\emph{输入退化时能否优雅降级}这两个问题却几乎无人系统回答。

围绕这一判断，我们给出三个模块：DAA 以低秩适配在不超过 $5$\,M 可训练参数的预算内缓解域偏移；EAD 把解释改造为带证据引用的可证伪断言，并以 InfoNCE 形式的证据对齐损失训练；DAG 以轻量质量探针驱动两条打分路径的连续切换，并在低置信时显式弃权。配套地，我们把解释可验证性（EAR / CFS / HR / TCR）与退化鲁棒性（RPR / ECE）提升为与检测精度并列的评测维度，并在 UCF-Crime 与三个跨域集、以及 $4\times3$ 退化网格上建立了统一协议。结果显示三者在互不重叠的维度上互补，且组合后存在弱的正向协同。

\paragraph{局限。} 第一，框架的全部结论目前建立在指标之上，尚未经过真实训练验证；第二，EAR / CFS 依赖「最小充分证据集 + 干扰证据 + 反事实标签」的人工标注，标注成本较高，其可扩展性有待检验；第三，DAG 的质量标签由退化类型与强度自动生成，对训练期未覆盖的新型退化可能存在分布漂移；第四，本文的评测仍在单视角数据上进行，跨摄像头视角一致性尚未纳入。

\paragraph{展望。} 未来工作有三个方向。其一是\textbf{跨数据集与跨视角泛化}：把 DAA 的「一场景一适配器」扩展为可组合的多场景适配器库，并在多摄像头场景中引入跨视角证据一致性约束，使结论不依赖单一数据集。其二是\textbf{实时推理与边缘部署}：结合边缘-云协同与知识图谱推理，把 DAG 的片段级动态路由扩展为设备级调度，在延迟预算内决定「哪些片段值得升级到大模型」。其三是\textbf{多模态扩展}：把音频、深度与热成像等模态纳入证据槽位体系，使 EAD 的「证据」不再局限于 RGB 时间区间，而成为跨模态、可核验的断言集合。我们相信，随着解释评测与鲁棒性评测成为社区共识，视频异常检测的评价体系将从「分数高低」走向「结论可信」。
